\subsection{Exponent Properties for Logarithms} 
The familiar rules for exponents can be rewritten as a rule for logarithms. In particular, we will be interested in the rules involving powers of the same \makeblank[1in]:
\begin{theorem}[Rules of Exponents]
For any real numbers $b$, $m$, $n$, and $p$ where $b>0$:
\begin{itemize}
\item $b^m\cdot b^n = b^{m+n}$ (product of powers rule)
\item $\dfrac{b^m}{b^n} = b^{m-n}$ (quotient of powers rule)
\item $\left(b^m\right)^p=b^{mp}$ (power of a power rule)
\end{itemize}
\end{theorem}
\begin{proof}
These rules are easy to show for \makeblank exponents. Then we defined all other exponents (negative integers, fractions, real numbers\textsuperscript{*}) to make the rules continue to hold.
\end{proof}
To get the same rules for logarithms, we will define:
\[
x=b^m
\qquad
y=b^n
\]
This means that $\log_b(x)=\makeblanksmall$ and $\log_b(y)=\makeblanksmall$.
\vfill
\textsuperscript{*} You might be thinking we didn't really define irrational exponents... and you're right. We use the fact that we can approximate any irrational number by a terminating decimal---which is rational---as closely as we'd like.